% Source: https://artofproblemsolving.com/community/c3h1908690_algebra
% Stable identity: aops-topic-1908690-q01
The expression is defined only for $x\ge0$. Put $t=\sqrt x\ge0$; then $x=t^2$ and $t$ runs over all nonnegative reals as $x$ does. For $t\ge0$,
\[
\sin t\ge0\iff 2k\pi\le t\le(2k+1)\pi\ \text{ for some integer }k\ge0 .
\]
Squaring (both ends are nonnegative, and squaring is increasing on $[0,\infty)$) gives
\[
\boxed{\,4k^2\pi^2\le x\le(2k+1)^2\pi^2\quad\text{for some integer }k\ge0\,}
\]
that is, $x\in[0,\pi^2]\cup[4\pi^2,9\pi^2]\cup[16\pi^2,25\pi^2]\cup\cdots$. The graph shows these intervals in red: the intervals where the sine is negative, $\left((2k+1)^2\pi^2,(2k+2)^2\pi^2\right)$, get longer as $x$ grows.

\begin{asy}
usepackage("amsmath");
size(9cm);
defaultpen(fontsize(10pt));
import graph;
real f(real t){return 1.4*sin(sqrt(20*t));}
pen hl=red+linewidth(3);
draw((0,0)--(0.4935,0),hl);draw((1.974,0)--(4.441,0),hl);draw((7.896,0)--(8.7,0),hl);
draw((0,0)--(9.2,0),Arrow(TeXHead));draw((0,-1.8)--(0,1.9),Arrow(TeXHead));
draw(graph(f,0,8.7,400),blue+linewidth(0.9));
draw((0,1.4)--(8.7,1.4),dotted);draw((0,-1.4)--(8.7,-1.4),dotted);
label("$1$",(0,1.4),W);label("$-1$",(0,-1.4),W);
label("$x$",(9.2,0),E);label("$y$",(0,1.9),N);
real[] xs={0.4935,1.974,4.441,7.896};
string[] ls={"$\pi^2$","$4\pi^2$","$9\pi^2$","$16\pi^2$"};
for(int i=0;i<4;++i){draw((xs[i],-0.1)--(xs[i],0.1));draw((xs[i],0)--(xs[i],-1.6),gray(0.5)+dotted);label(ls[i],(xs[i],-1.6),S,fontsize(9pt));}
label("$y=\sin\sqrt x$",(3.2,1.45),N,blue);
label("\mbox{}",(min(currentpicture,user=true)+max(currentpicture,user=true))/2,black); // keeps the text after the figure black when LaTeX typesets the labels inline
\end{asy}
